\subsection{Uniqueness Theorem for Homogeneous Means} \label{sect-HLP_proof}
In this section we will prove the uniqueness theorem for homogeneous means, which states that the power means \footnote{In some literature this is also referred to as the Hölder mean or the generalized means.} (including the geometric mean as a limit case) are the only homogeneous quasi-arithmetic means.
Before beginning the proof, we will first define quasi-arithmetic means and one class of this means, the power means.
Both, the proof and the definitions are based on \parencite[pp. 66-69]{hardy1952inequalities} and \parencite[p. 271-273]{handbookMeans}.

We begin with defining quasi-arithmetic means. The name stems from generalizing the arithmetic mean by applying a function on each value before averaging and then using its inverse to transform the result back.

\begin{defi}[Quasi Arithmetic Means and Homogeneity] \index{Quasi-arithmetic mean} \index{Quasi-arithmetic mean!homogeneous} %\gls{quasi-arith-mean}
	Let $\phi : [0,\infty) \rightarrow [0,\infty)$ be continuous and strictly increasing and let $\ve{x}, \ve{w} \in \mathbb{R}_{\geq 0}^n$.
	Then the quasi-arithmetic $\phi$-mean of $\ve{x}$ with weights $\ve{w}$
	$$\mean{\ve{x}} = \phi^{-1} \left ( \sum_{i=1}^{n} w_i \phi(x_i) \right)$$
	The mean is said to be \emph{homogeneous}, if
	$$\mean{t \cdot \ve{x}} = t \cdot \mean{\ve{x}}$$
	holds for all $t > 0$.
\end{defi}

Next, we define the power means and remark that they are quasi-arithmetic means.

\begin{defi}[Power mean] \index{Power mean} %\gls{powermean}
	For $n \in \mathbb{N}, \ve{x}, \ve{w} \in \mathbb{R}_{\geq 0}^n$ and $ p > 0$ we define the power means $\meanp{\ve{x}}$ as
	$$ \meanp{\ve{x}} = \left ( \frac{1}{\sum_{i=1}^{n} w_i} \sum_{i=1}^{n}w_i x_i^p \right)^{1/p}. $$
	For $p=0$ we can define $\meanp[0]{\ve{x}}$ as the geometric mean
	$$
\meanp[0]{\ve{x}} = \left ( \prod_{i=1}^n x_i^{w_i} \right )^{1 / \sum_{i=1}^{n}w_i}
	$$
	since it is obtained as a limit case for $r \rightarrow 0$ \parencite[p. 176]{handbookMeans}.
\end{defi}

\begin{rem}[Power Means are Quasi-Arithmetic]
	By setting $\phi(x) = x^p$ for $p>0$ or $\phi(x) = \log(x)$ we obtain the power means for $p>0$ and $p=0$ respectively.\\
  In particular, setting $w_i = 1$ yields the familiar generalizations of the power mean
	$$\left ( \frac{1}{n} \sum_{i=1}^{n}x_i^p \right)^{1/p}. $$
  and the arithmetic mean
  $$\sqrt[n]{\prod_{i=1}^n x_i}.$$
\end{rem}

Now we can continue with a theorem about equal means, which is needed in the proof of the uniqueness theorem for homogeneous means.

\begin{thm}[Characterization of Equal Means] \label{thm-quasi-mean-equal-cond}
	In order that
  \begin{equation} \label{thm-quasi-mean-eq-equal}
    \mean[\chi]{\ve{a}}  = \mean[\psi]{\ve{a}}
  \end{equation}
	for all $\ve{a}, \ve{w}$ and $n \in \mathbb{N}$, it is necessary and sufficient that
  \begin{equation} \label{thm-quasi-mean-eq-linear}
		\phi = \alpha \cdot  \psi + \beta  \end{equation}
	where $\alpha$ and $\beta$ are constants and $\alpha \neq 0$.
\end{thm}

\begin{proof}

  \begin{proofpart}[$\Rightarrow$]
    Assume that \Cref{thm-quasi-mean-eq-equal} holds. In particular, it also holds for $n=2$.
    Thus we have
    \begin{equation*}
      \chi^{-1}(w_1 \chi(a_1) + w_2 \chi(a_2)) = \psi^{-1} (w_1 \psi(a_1) + w_2 \psi (a_2))
    \end{equation*}

    Putting $\phi(x) = \chi ( \psi^{-1} (x)), x_i = \psi (a_i)$ for $i=1, 2$ yields
    \begin{equation*}
      \phi(w_1 x_1 + w_2 x_2) = w_1 \psi(x_1) + w_2 \psi(x_2) \quad x_1, x_2 \in \mathbb{R}_{\geq  0}.
    \end{equation*}
    In particular, this holds for $w_1, w_2 = 1/2$ which then reduces to
    \begin{equation*}
      \phi \left (\frac{x_1 + x_2}{2}\right ) = \frac{\phi(x_1)}{2} + \frac{\phi (x_2)}{2} \quad x_1, x_2 \in \mathbb{R}_{\geq  0}.
    \end{equation*}

    By \parencite[p. 46]{aczel1966lectures}, the general continuous solution is $\psi(x) = \alpha x + \beta$, hence $\chi = \alpha \psi + \beta$.
  \end{proofpart}

  \begin{proofpart}[$\Leftarrow$]
    Assume that \Cref{thm-quasi-mean-eq-linear} holds.
    Then
    \begin{align*}
      \chi ( \mean[\chi]{\ve{a}}) &= \sum_{i=1}^{n} w_i \chi(a_i) \tag{cancellation of inverses} \\
                                  &= \sum_{i=1}^{n} q_i \alpha \psi(a_i) + \beta \tag{\cref{thm-quasi-mean-eq-linear}} \\
                                  &= \alpha \left ( \sum_{i=1}^{n} w_i \psi(a_i) \right ) + \beta \tag{rearranging} \\
                                  &= \chi \left [ \mean[\psi]{\ve{a}} \right] \tag{\cref{thm-quasi-mean-eq-linear}}
    \end{align*}
    Applying $\chi^{-1}$ to both sides of the equation yields \Cref{thm-quasi-mean-eq-equal}.
  \end{proofpart}
\end{proof}

Now we finally state and prove the main theorem of this section.

\begin{thm}[$\meanpempty$ are the only homogeneous means $\meanempty$] \label{thm-power-mean-uniqueness}
	Suppose that $\phi$ is continuous and strictly increasing in the open interval $(0,\infty)$ and that the mean defined by $\phi$ is homogeneous, i.e.
  \begin{equation} \label{thm-quasi-homogen-eq}
		\mean{t \ve{a} } = t \mean{\ve{a}}
  \end{equation}
	for all positive $\ve{a}, \ve{w},t$ and all $n \in \mathbb{N}$ such that $\sum_{i=1}^{n} w_i = 1$. Then $\mean{ \ve{a}}$ is $\meanp{\ve{a} }$.

\end{thm}

\begin{proof}
  By \cref{thm-quasi-mean-equal-cond} we can assume that $\phi(1) = 0$, since we can replace $\phi(x)$ by $\phi(x) - \phi(1)$ without changing $\mean{\ve{a}}$. Next, we write \Cref{thm-quasi-homogen-eq} as
  \begin{align*}
    \mean{\ve{a}} = t^{-1} \mean{t\ve{a}} = t^{-1} \phi^{-1} \left ( \sum_{i=1}^{n} w_i \phi(t a_i) \right) = \mean[\psi]{\ve{a}}
  \end{align*}
  by setting $\psi(x) = \phi(t x)$. From this, it follows that $\psi^{-1}(x) = t^{-1} \cdot \phi(x)$ \footnote{For this, let $\psi^{-1}(x) = y$. Then, by applying $\psi$ to both sides, and using its definition, this yields $x = \phi(ty)$. Applying the inverse of $\phi$ and dividing by $t$ yields $t^{-1} \phi^{-1}(x) = y$, which in turn is equal to $\psi^{-1}(x)$ by the first equation.}, which in turn justifies the above equation.
  Therefore, we have the equality of the means defined by $\psi$ and $\phi$. Since $\psi$ is defined for every $t > 0$, we have by \cref{thm-quasi-mean-equal-cond} that there are $\alpha(t), \beta(t) \in \mathbb{R}, \alpha(t) \neq  0$ such that
  \begin{equation*}
  \phi(t x) = \psi(x) = \alpha(t) \cdot \phi(x) + \beta(t).
  \end{equation*}
  Plugging in $x=1$ and using $\phi(1) = \phi(0)$ yields
  \begin{equation*}
  \phi(t) = \psi(1) = \beta(t).
  \end{equation*}
  Thus, for all positive $x, y$ by setting $y=t$
  \begin{equation} \label{proof-quasi-eq2}
    \phi(x y) = \alpha(y)\phi(x) + \phi(y)
  \end{equation}
  Similarly, by exchanging the roles of $x$ and $y$ we get
  \begin{equation*}
    \phi(x y) = \alpha(x)\phi(y) + \phi(x)
  \end{equation*}
  Combining these two equalities yields
  \begin{align*}
    \alpha(y)\phi(x) + \phi(y)  &= \alpha(x)\phi(y) + \phi(x) \tag{$- \phi(x) - \phi(y)$} \\
    \alpha(y)\phi(x) - \phi(x)  &= \alpha(x)\phi(y) - \phi(y) \tag{divide \footnotemark by $\phi(x) \cdot \phi(y)$} \\
    \frac{\alpha(x) -1}{\phi(x)} &= \frac{\alpha(y) - 1}{\phi(y)} := c
  \end{align*} \footnotetext{Provided $x \neq 1, y \neq 1$. Since \Cref{proof-quasi-eq1} is true when $x$ or $y$ is $1$, we can ignore the exception.}
  Since this must hold for all $x$ and $y$, the expressions must reduce to a constant $c$. This can be seen by settings, for example, $x=1$ and observing that the right hand side must equal to that for all $y$ and therefore must be constant. The same argument can be applied in reverse.
  Using this, we get the following expression
  \begin{equation*}
    \alpha(y) = 1 + c \phi(y).
  \end{equation*}
  Plugging this expression for $\alpha(y)$ into \Cref{proof-quasi-eq2} yields
  \begin{equation} \label{proof-quasi-eq1}
    \phi(xy) = c \phi(x) \phi(y) + \phi(x) + \phi(y).
  \end{equation}

  For this equation, there are two cases to consider: $c=0$ and $c \neq 0$.
  If $c=0$, then \Cref{proof-quasi-eq1} reduces to the well known logarithmic Cauchy equation
  $$  \phi(xy) = \phi(x) + \phi(y), $$
  of which the most general and for $x>0$ continuous solution is $\phi(x) = \alpha \log(x)$ for some constant $\alpha \in \mathbb{R}$; see \parencite[p. 39]{aczel1966lectures} and \parencite[p. 110-111]{cauchy1821}.

  If $c \neq 0$, we set $f(x) = c \phi(x) + 1$ and \Cref{proof-quasi-eq1} reduces to the multiplicative Cauchy functional equation
  $$ f(xy) = f(x)f(y), $$
  whose general solution is $f(x) = x^r$; see \parencite[p. 39]{aczel1966lectures} and \parencite[p. 111-112]{cauchy1821}. Substituting back yields $\phi(x) = \frac{x^r -1}{c}$.
\end{proof}

\begin{sourcescontribs}
  The main results of this chapter, namely \cref{thm-conds} and \cref{thm-uniquene-repr-minkowski}, were first published in \parencite{foundationsMDS}, then proved in \parencite{tversky1970dimensional} and finally refined and updated to the version presented in \parencite[Chapter 14.4.5, pp. 185-187]{foundations2}.
  This chapter is primarily based on the latter.\\
  The proofs presented in \cref{sec-conds} extend some arguments from \parencite[pp. 192-194]{foundations2} and include additional references.\\
  %

  The proofs presented in \cref{sec-unique-repr-thm} largely extend and complete the proofs provided in \parencite[pp. 195-199]{foundations2} and \parencite[p. 589-593]{tversky1970dimensional} \footnote{Until the end of this section, all references are with respect to these two sources, e.g., when writing "the proofs are briefly outlined"}. In both sources, the statements and proofs are almost identical.
  The proof of the main theorem, \cref{thm-uniquene-repr-minkowski}, can mostly be found in both sources.
  However, important details are missing in this proof. Most notably, the preconditions for the application of \cref{thm-power-mean-uniqueness} have neither been mentioned nor established.

  The proofs of the homogeneity theorem (\cref{sect-homogeneity-theorem}) and the convexity of $\metric$ balls (\cref{sec-ball-convexity}) are briefly outlined; however, some arguments are either skipped or only briefly mentioned without further explanation.
  Moreover, the statement of the homogeneity theorem is incorrect, as continuity is used but not assumed.
  \cref{lemma-part1} and \cref{lemma-part2} of the proof can be found in both references. The other parts of the proof are briefly outlined, so the majority of the work in this section involves the author completing the proofs.
  \cref{lemma-distance} was inspired by a post on math stackexchange, see \href{https://math.stackexchange.com/questions/48714/a-and-b-disjoint-a-compact-and-b-closed-implies-there-is-positive-dist}{https://math.stackexchange.com/questions/48714/a-and-b-disjoint-a-compact-and-b-closed-implies-there-is-positive-dist}.\\
  The definitions and convexity results in \cref{sec-ball-convexity} are based on Chapters 1.1 to 1.4 of \parencite[pp. 1-38]{lecturesConvGeo}. Most other arguments have been briefly outlined.

  The section on homogeneous quasi-arithmetic means, \cref{sect-HLP_proof}, is based on the work found in \parencite[pp. 66-69]{hardy1952inequalities} and \parencite[p. 271-273]{handbookMeans}. The proofs provided here can be found in these references.

\end{sourcescontribs}

